\documentclass[11pt,openany]{report}
\usepackage[letterpaper,margin=0.82in,headheight=14pt]{geometry}
\usepackage[T1]{fontenc}
\usepackage[utf8]{inputenc}
\usepackage{lmodern}
\usepackage{microtype}
\usepackage{amsmath,amssymb,amsthm,mathtools}
\usepackage{booktabs,longtable,tabularx,array,multirow}
\usepackage{graphicx,float,placeins}
\usepackage[dvipsnames]{xcolor}
\usepackage{enumitem}
\usepackage{fancyhdr}
\usepackage{listings}
\usepackage{url}
\usepackage[hidelinks]{hyperref}
\hypersetup{
  pdftitle={Projector-Generated Special Functions I},
  pdfauthor={Komorebi P.},
  pdfsubject={Polyphase differential lifts, Bernoulli-Lerch wavelets, cyclotomic details, and Mobius coherence}
}
\usepackage[nameinlink,noabbrev]{cleveref}
\usepackage{needspace}
\allowdisplaybreaks

\definecolor{Ink}{HTML}{17202A}
\definecolor{Blue}{HTML}{2E86AB}
\definecolor{Teal}{HTML}{2A9D8F}
\definecolor{Orange}{HTML}{F4A261}
\definecolor{Magenta}{HTML}{A23E8C}
\definecolor{Red}{HTML}{E76F51}
\definecolor{Paper}{HTML}{FAFAF7}
\definecolor{Muted}{HTML}{5D6D7E}
\definecolor{CodeBG}{HTML}{F3F5F5}

\pagecolor{Paper}
\color{Ink}
\graphicspath{{figures/}}
\setlength{\parindent}{0pt}
\setlength{\parskip}{0.52em}
\raggedbottom

\pagestyle{fancy}
\fancyhf{}
\fancyhead[L]{\small Projector-Generated Special Functions I}
\fancyhead[R]{\small Komorebi P.}
\fancyfoot[C]{\thepage}
\renewcommand{\headrulewidth}{0.4pt}

\newtheorem{theorem}{Theorem}[chapter]
\newtheorem{proposition}[theorem]{Proposition}
\newtheorem{lemma}[theorem]{Lemma}
\newtheorem{corollary}[theorem]{Corollary}
\newtheorem{conjecture}[theorem]{Conjecture}
\theoremstyle{definition}
\newtheorem{definition}[theorem]{Definition}
\newtheorem{example}[theorem]{Worked example}
\newtheorem{algorithm}[theorem]{Algorithm}
\theoremstyle{remark}
\newtheorem{remark}[theorem]{Remark}
\newtheorem{boundary}[theorem]{Boundary statement}

\newcommand{\C}{\mathbb C}
\newcommand{\R}{\mathbb R}
\newcommand{\N}{\mathbb N}
\newcommand{\Z}{\mathbb Z}
\newcommand{\Q}{\mathbb Q}
\newcommand{\e}{\mathrm e}
\newcommand{\ii}{\mathrm i}
\newcommand{\Bfrak}{\mathfrak B}
\newcommand{\Wcal}{\mathcal W}
\newcommand{\Ccal}{\mathcal C}
\newcommand{\Rcal}{\mathcal R}
\newcommand{\Mcal}{\mathcal M}
\newcommand{\norm}[1]{\left\lVert#1\right\rVert}
\newcommand{\abs}[1]{\left\lvert#1\right\rvert}
\newcommand{\ip}[2]{\left\langle#1,#2\right\rangle}
\newcommand{\dd}{\,\mathrm d}
\newcommand{\Res}{\operatorname{Res}}
\newcommand{\diag}{\operatorname{diag}}
\newcommand{\Span}{\operatorname{span}}
\newcommand{\rank}{\operatorname{rank}}

\newcolumntype{L}[1]{>{\raggedright\arraybackslash}p{#1}}

\lstdefinestyle{pythonstyle}{
  language=Python,
  basicstyle=\ttfamily\scriptsize,
  backgroundcolor=\color{CodeBG},
  frame=single,
  rulecolor=\color{Muted!45},
  breaklines=true,
  breakatwhitespace=false,
  showstringspaces=false,
  numbers=left,
  numberstyle=\tiny\color{Muted},
  numbersep=7pt,
  tabsize=4,
  columns=fullflexible,
  keepspaces=true
}
\lstset{style=pythonstyle}

\newcommand{\FigureWide}[3]{%
\Needspace{0.53\textheight}
\begin{figure}[H]
\centering
\includegraphics[width=0.92\textwidth,height=0.44\textheight,keepaspectratio]{#1.pdf}
\caption{#2 #3}
\label{fig:#1}
\end{figure}
}

\newcommand{\InlinePart}[1]{%
\FloatBarrier
\Needspace{10\baselineskip}
\refstepcounter{part}%
\addcontentsline{toc}{part}{\protect\numberline{\thepart}#1}%
\vspace{1.4em}
{\color{Blue}\Large\bfseries Part \thepart\par}
\vspace{0.2em}
{\color{Blue}\LARGE\bfseries #1\par}
\vspace{0.9em}
}

\newcommand{\FlowChapter}[1]{%
\FloatBarrier
\Needspace{14\baselineskip}
\refstepcounter{chapter}%
\setcounter{section}{0}%
\setcounter{subsection}{0}%
\addcontentsline{toc}{chapter}{\protect\numberline{\thechapter}#1}%
\markboth{#1}{#1}%
\vspace{1.3em}
{\color{Muted}\large\bfseries Chapter \thechapter\par}
\vspace{0.15em}
{\LARGE\bfseries #1\par}
\vspace{0.9em}
}

\title{\textbf{Projector-Generated Special Functions I}\\[0.25em]
\Large Polyphase Differential Lifts, Bernoulli-Lerch Wavelets, Cyclotomic Detail Functions, and M\"obius Coherence}
\author{Komorebi P.}
\date{}

\begin{document}
\maketitle

\begin{abstract}
This paper develops a special-function theory generated by exact projector and polyphase operations. For an analytic scalar function, residue sections are assembled into a vector-valued packet function. Multiplication and the Euler derivative become a carry-shift matrix and a diagonal differential operator that satisfy the same Weyl commutator as the scalar pair. Every ordered scalar differential equation therefore lifts to an exact matrix differential system.

The construction is applied to exponentials, Bessel functions, Airy functions, generalized hypergeometric equations, Bernoulli interpolation, Hurwitz zeta, Lerch transcendent, and M\"obius-weighted reciprocal $L$-function transfers. A complex-order Bernoulli-Hurwitz packet family is proved to obey exact reconstruction, exact parent-child refinement, an order-lowering differential ladder, a translation law, and integer interpolation of the terms in Raabe multiplication. A unitary transform across the child address yields cyclotomic Bernoulli-Lerch details with an inverse transform and Parseval identity. The two-channel case identifies Bernoulli and Euler polynomials as the low-pass and high-pass outputs of one exact refinement bank.

Packet Wronskians provide a direct test for ratio monotonicity, zero interlacing, and Hermite-Biehler compatibility. Computation also supplies a useful negative result: the first Bernoulli-Euler Wronskian family is not universally real-rooted. The final part introduces a generalized M\"obius-Bernoulli transfer function. Its binomial, Dirichlet, exponential-generating, Mellin, wavelet, polyphase, character, and quantum forms are proved. The diagonal packet energy has a closed asymptotic law. In the zeta case, the Riemann hypothesis is equivalent to subpolynomial growth of a normalized packet-coherence ratio. This equivalence is a theorem target, not a proof of the hypothesis.

All figures were generated from the accompanying Python program. The release contains 26 vector figures, matching PNG files, exact tables, 51 algebraic or high-precision checks, and the full source listing.
\end{abstract}

\chapter*{Status and claim boundary}
\addcontentsline{toc}{chapter}{Status and claim boundary}
The paper proves the algebraic lift, the Bernoulli-Hurwitz packet laws, the cyclotomic transform, the Bernoulli-Euler channel identity, the packet-Wronskian formulas, the generalized transfer identities, the Mellin formula, the long-time power identity, and the packet-energy asymptotic under the stated hypotheses.

Classical inputs include the Hurwitz zeta function, the Lerch transcendent, Raabe multiplication, Bernoulli and Euler polynomial identities, inverse Euler products, the Riesz and B\'aez-Duarte criteria, and standard squarefree-density estimates.

The paper does not claim that every generated function is historically new. The contribution is the common reconstruction, refinement, differential, transform, and certificate structure. The Bernoulli interpolation and its distinction from the integer Bernoulli-value sequence were developed in the preceding transfer-block note~\cite{KomorebiBernoulli}. Historical novelty of individual specializations requires a separate literature study. The packet-coherence condition in the zeta case is equivalent to RH and remains open.

\tableofcontents
\listoffigures

\InlinePart{The general polyphase lift}

\FlowChapter{From a scalar germ to a vector special function}

Let
\[
F(z)=\sum_{n=0}^{\infty}a_nz^n
\]
be analytic in a neighborhood of the origin. Fix an integer $M\ge2$.

\begin{definition}[Residue-section functions]
For $0\le r<M$, define
\[
\boxed{
F_{M,r}(u)=\sum_{m=0}^{\infty}a_{Mm+r}u^m.
}
\]
The packet vector is
\[
\mathbf F_M(u)=
\begin{pmatrix}
F_{M,0}(u)\\
F_{M,1}(u)\\
\vdots\\
F_{M,M-1}(u)
\end{pmatrix}.
\]
\end{definition}

\begin{theorem}[Exact polyphase reconstruction]
Wherever the defining series converge,
\[
\boxed{
F(z)=\sum_{r=0}^{M-1}z^rF_{M,r}(z^M).
}
\]
\end{theorem}

\begin{proof}
Every exponent has a unique representation $n=Mm+r$ with $0\le r<M$. Therefore
\[
\sum_{n\ge0}a_nz^n
=
\sum_{r=0}^{M-1}
\sum_{m\ge0}a_{Mm+r}z^{Mm+r}
=
\sum_{r=0}^{M-1}z^rF_{M,r}(z^M).
\]
\end{proof}

\FigureWide{01_polyphase_partition}{A coefficient sequence separated into four residue packets.}{Each Taylor exponent appears in exactly one row, so the analysis map loses no coefficient.}

This elementary identity becomes more useful when the scalar differential operators are also transported into packet coordinates.

\FlowChapter{The carry-shift representation of the Weyl algebra}

Define the carry-shift matrix
\[
C_M(u)=
\begin{pmatrix}
0&0&\cdots&0&u\\
1&0&\cdots&0&0\\
0&1&\cdots&0&0\\
\vdots&&\ddots&&\vdots\\
0&0&\cdots&1&0
\end{pmatrix}.
\]
Define the packet Euler operator
\[
\Theta_M
=
Mu\frac{\dd}{\dd u}I_M+
\diag(0,1,\ldots,M-1).
\]

\begin{proposition}[Intertwining laws]
Let
\[
\mathcal S_M\mathbf f(z)
=
\sum_{r=0}^{M-1}z^rf_r(z^M).
\]
Then
\[
\boxed{
\mathcal S_M(C_M\mathbf f)=z\,\mathcal S_M\mathbf f
}
\]
and
\[
\boxed{
\mathcal S_M(\Theta_M\mathbf f)
=
\left(z\frac{\dd}{\dd z}\right)\mathcal S_M\mathbf f.
}
\]
\end{proposition}

\begin{proof}
Multiplication by $z$ sends packet $r$ to packet $r+1$. The last packet wraps to packet $0$ and contributes a factor $z^M=u$. This is exactly $C_M(u)$. For the differential identity, a monomial $u^m$ in packet $r$ reconstructs to $z^{Mm+r}$. The operator $Mu\frac{\dd}{\dd u}+r$ multiplies that monomial by $Mm+r$, which is the eigenvalue of $z\frac{\dd}{\dd z}$ on the reconstructed monomial.
\end{proof}

\begin{corollary}[Weyl commutator]
\[
\boxed{
[\Theta_M,C_M]=C_M.
}
\]
\end{corollary}

\begin{proof}
The scalar operators satisfy
\[
\left[z\frac{\dd}{\dd z},z\right]=z.
\]
Apply the two intertwining laws. The synthesis map is injective because the residue decomposition is unique.
\end{proof}

\FigureWide{02_carry_shift}{The carry-shift action for five packets.}{Multiplication advances the packet address, while the last packet wraps around with the factor $u$.}

\begin{theorem}[Differential lift]
Let $\mathcal L(z,\theta)$ be an ordered polynomial in the noncommuting variables $z$ and $\theta$, where $\theta=z\frac{\dd}{\dd z}$. If
\[
\mathcal L(z,\theta)F=0,
\]
then
\[
\boxed{
\mathcal L(C_M,\Theta_M)\mathbf F_M=0.
}
\]
The order of the factors in every monomial of $\mathcal L$ is preserved.
\end{theorem}

\begin{proof}
The assignment
\[
z\longmapsto C_M,
\qquad
\theta\longmapsto\Theta_M
\]
preserves the defining commutator. It therefore gives a representation of the same ordered operator algebra. Intertwining each ordered monomial with $\mathcal S_M$ gives
\[
\mathcal S_M\bigl(\mathcal L(C_M,\Theta_M)\mathbf F_M\bigr)
=
\mathcal L(z,\theta)F=0.
\]
Injectivity of $\mathcal S_M$ gives the result.
\end{proof}

\FigureWide{05_operator_lift_map}{The operator-level lift.}{The scalar differential equation and the packet matrix system are two coordinate realizations of the same ordered algebra.}

\FlowChapter{Worked lifts of classical functions}

\section{The exponential}
The scalar equation
\[
\theta e^z=ze^z
\]
lifts to
\[
\boxed{
\Theta_M\mathbf E_M=C_M\mathbf E_M.
}
\]
The residue sections are
\[
E_{M,r}(u)=\sum_{m=0}^{\infty}\frac{u^m}{(Mm+r)!}.
\]
Gauss multiplication gives
\[
\boxed{
E_{M,r}(u)
=
\frac1{r!}
{}_0F_{M-1}
\left(
\begin{matrix}-\\
\frac{r+1}{M},\ldots,\widehat{1},\ldots,\frac{r+M}{M}
\end{matrix}
\middle|
\frac{u}{M^M}
\right),
}
\]
where the denominator parameter equal to $1$ is omitted.

\begin{example}[Two channels]
For $M=2$,
\[
E_{2,0}(u)=\cosh\sqrt u,
\qquad
E_{2,1}(u)=\frac{\sinh\sqrt u}{\sqrt u}.
\]
Hence
\[
e^z=E_{2,0}(z^2)+zE_{2,1}(z^2)=\cosh z+\sinh z.
\]
\end{example}

\begin{example}[Three channels]
For $M=3$,
\[
\begin{aligned}
E_{3,0}(u)
&={}_0F_2\left(;\frac13,\frac23;\frac{u}{27}\right),\\
E_{3,1}(u)
&={}_0F_2\left(;\frac23,\frac43;\frac{u}{27}\right),\\
E_{3,2}(u)
&=\frac12{}_0F_2\left(;\frac43,\frac53;\frac{u}{27}\right).
\end{aligned}
\]
The exact synthesis is
\[
\boxed{
e^z=E_{3,0}(z^3)+zE_{3,1}(z^3)+z^2E_{3,2}(z^3).
}
\]
\end{example}

\FigureWide{03_exponential_sections}{Exponential packets for two and three channels.}{The section functions are ordinary generalized hypergeometric functions arranged as one exact vector object.}
\FigureWide{04_exponential_reconstruction}{Numerical reconstruction residuals for three packet counts.}{The residual remains at floating-point roundoff across the displayed interval.}

\section{Bessel functions}
The Bessel equation for $J_\nu$ can be written
\[
\left(\theta^2+z^2-\nu^2\right)J_\nu(z)=0.
\]
Whenever the chosen branch is analytic in the working neighborhood, the packet vector satisfies
\[
\boxed{
\left(\Theta_M^2+C_M^2-\nu^2I\right)\mathbf J_{\nu,M}=0.
}
\]
For integer order, the Taylor expansion at the origin supplies a global entire packet vector. Figure~\ref{fig:06_bessel_sections} shows the three residue sections of $J_0$.

\FigureWide{06_bessel_sections}{Three residue sections of $J_0$.}{The scalar oscillation is divided by exponent class, while exact synthesis recovers the original Bessel function.}

\section{Airy functions}
The Airy equation $F''=zF$ is equivalent to
\[
\theta(\theta-1)F=z^3F.
\]
Its packet lift is
\[
\boxed{
\Theta_M(\Theta_M-I)\mathbf F_M=C_M^3\mathbf F_M.
}
\]
This example is useful because the scalar equation contains a derivative order and a multiplication order that do not commute. The ordered lift preserves both.

\section{Generalized hypergeometric functions}
The generalized hypergeometric function satisfies
\[
\left[
\theta\prod_{j=1}^{q}(\theta+b_j-1)
-z\prod_{i=1}^{p}(\theta+a_i)
\right]F=0.
\]
The packet system is therefore
\[
\boxed{
\left[
\Theta_M\prod_{j=1}^{q}(\Theta_M+(b_j-1)I)
-C_M\prod_{i=1}^{p}(\Theta_M+a_iI)
\right]\mathbf F_M=0.
}
\]
This gives a systematic matrix-valued lift of every scalar generalized hypergeometric equation.

\InlinePart{Bernoulli-Hurwitz packets and cyclotomic details}

\FlowChapter{Complex-order Bernoulli interpolation}

For $a>0$, define
\[
\boxed{
\Bfrak(z;a)=-z\zeta(1-z,a).
}
\]
The pole of the Hurwitz zeta function at $1-z=1$ is removed by the factor $-z$, so $\Bfrak(z;a)$ is entire in $z$ and $\Bfrak(0;a)=1$. The special-value and multiplication formulas used below are classical properties of Bernoulli polynomials and the Hurwitz zeta function~\cite{DLMF24,DLMF25}.

For positive integers $n$,
\[
\Bfrak(n;a)=B_n(a).
\]
Thus $\Bfrak$ is a complex-order interpolation of the Bernoulli polynomial sequence.

\begin{definition}[Bernoulli-Hurwitz packet]
For $M\ge1$ and $0\le r<M$, define
\[
\boxed{
\Bfrak_{M,r}(z;a)
=-zM^{z-1}
\zeta\left(1-z,\frac{a+r}{M}\right).
}
\]
At $z=0$, define $\Bfrak_{M,r}(0;a)=1/M$ by removal of the singularity.
\end{definition}

\begin{theorem}[Reconstruction]
\[
\boxed{
\Bfrak(z;a)=\sum_{r=0}^{M-1}\Bfrak_{M,r}(z;a).
}
\]
\end{theorem}

\begin{proof}
The Hurwitz multiplication formula is
\[
\sum_{r=0}^{M-1}
\zeta\left(s,\frac{a+r}{M}\right)
=M^s\zeta(s,a).
\]
Set $s=1-z$ and multiply by $-zM^{z-1}$.
\end{proof}

\FigureWide{07_bernoulli_packet_reconstruction}{Four Bernoulli-Hurwitz packets and their synthesis.}{The packets may be individually large or oscillatory, but their complete sum equals the direct interpolation.}

\FlowChapter{Exact refinement, differentiation, and translation}

\begin{theorem}[Parent-child refinement]
For integers $L,M\ge1$,
\[
\boxed{
\Bfrak_{M,r}(z;a)
=
\sum_{\ell=0}^{L-1}
\Bfrak_{LM,r+\ell M}(z;a).
}
\]
\end{theorem}

\begin{proof}
Set $A=(a+r)/M$. The child arguments are $(A+\ell)/L$. Summing their Hurwitz zeta functions and applying the multiplication formula at modulus $L$ gives the parent. The powers of $L$ cancel exactly.
\end{proof}

\begin{theorem}[Integer samples]
For every positive integer $n$,
\[
\boxed{
\Bfrak_{M,r}(n;a)
=M^{n-1}
B_n\left(\frac{a+r}{M}\right).
}
\]
\end{theorem}

\begin{proof}
Use $\zeta(1-n,x)=-B_n(x)/n$.
\end{proof}

The reconstruction theorem at integer order becomes Raabe multiplication:
\[
B_n(a)=M^{n-1}\sum_{r=0}^{M-1}B_n\left(\frac{a+r}{M}\right).
\]

\begin{theorem}[Differential ladder]
\[
\boxed{
\frac{\partial}{\partial a}\Bfrak_{M,r}(z;a)
=z\Bfrak_{M,r}(z-1;a).
}
\]
\end{theorem}

\begin{proof}
The Hurwitz derivative is
\[
\frac{\partial}{\partial x}\zeta(s,x)=-s\zeta(s+1,x).
\]
Apply the chain rule to $x=(a+r)/M$ and simplify.
\end{proof}

\begin{theorem}[Translation law]
\[
\boxed{
\Bfrak_{M,r}(z;a+M)-\Bfrak_{M,r}(z;a)
=z(a+r)^{z-1}.
}
\]
\end{theorem}

\begin{proof}
Use
\[
\zeta(s,x+1)=\zeta(s,x)-x^{-s}.
\]
\end{proof}

\FigureWide{08_bernoulli_refinement_tree}{A Bernoulli packet refined from modulus one to modulus eight.}{Every parent is the exact sum of its children. The refinement law works at complex order, not only at integer samples.}

\FlowChapter{Cyclotomic Bernoulli-Lerch details}

Let $\omega_L=\e^{2\pi\ii/L}$.

\begin{definition}[Cyclotomic detail]
\[
\boxed{
\Wcal_{M,r}^{(L,j)}(z;a)
=
\frac1{\sqrt L}
\sum_{\ell=0}^{L-1}
\omega_L^{-j\ell}
\Bfrak_{LM,r+\ell M}(z;a).
}
\]
\end{definition}

\begin{theorem}[Lerch closed form]
\[
\boxed{
\Wcal_{M,r}^{(L,j)}(z;a)
=
-\frac{z}{\sqrt L}
M^{z-1}
\Phi\left(
\omega_L^{-j},1-z,\frac{a+r}{M}
\right).
}
\]
\end{theorem}

\begin{proof}
Write each child Hurwitz zeta function as a series. The child address $\ell$ and the internal summation index combine into one integer $n=Lm+\ell$. The Fourier weight $\omega_L^{-j\ell}$ becomes $\omega_L^{-jn}$. The resulting series is the Lerch transcendent.
\end{proof}

The DC channel is the parent:
\[
\boxed{
\Wcal_{M,r}^{(L,0)}=\frac1{\sqrt L}\Bfrak_{M,r}.
}
\]
The inverse transform is
\[
\boxed{
\Bfrak_{LM,r+\ell M}
=
\frac1{\sqrt L}
\sum_{j=0}^{L-1}
\omega_L^{j\ell}
\Wcal_{M,r}^{(L,j)}.
}
\]
For each fixed $(z,a)$, Parseval gives
\[
\boxed{
\sum_{\ell=0}^{L-1}
\abs{\Bfrak_{LM,r+\ell M}}^2
=
\sum_{j=0}^{L-1}
\abs{\Wcal_{M,r}^{(L,j)}}^2.
}
\]

\FigureWide{09_cyclotomic_details}{Four cyclotomic Bernoulli-Lerch channels as the translation parameter varies.}{The non-DC details carry phase and magnitude information that disappears in the parent sum.}
\FigureWide{10_cyclotomic_parseval}{Energy before and after the cyclotomic transform.}{The child and detail energies agree to numerical precision, so the transform is a genuine perfect-reconstruction bank.}

\FlowChapter{Exponential generating function of the detail bank}

At positive integer order, define
\[
W_{L,j}(n;a)
=
\frac1{\sqrt L}
\sum_{\ell=0}^{L-1}
\omega_L^{-j\ell}
L^{n-1}B_n\left(\frac{a+\ell}{L}\right).
\]

\begin{theorem}[Cyclotomic EGF]
\[
\boxed{
\sum_{n=0}^{\infty}
W_{L,j}(n;a)\frac{t^n}{n!}
=
\frac{t\e^{at}}
{\sqrt L\left(\omega_L^{-j}\e^t-1\right)}.
}
\]
\end{theorem}

\begin{proof}
The EGF of the $\ell$-th child is
\[
\sum_{n\ge0}
L^{n-1}B_n\left(\frac{a+\ell}{L}\right)
\frac{t^n}{n!}
=
\frac{t\e^{(a+\ell)t}}{\e^{Lt}-1}.
\]
Take the finite Fourier transform over $\ell$ and sum the geometric progression. The factor $1-\e^{Lt}$ cancels the denominator.
\end{proof}

This is the standard root-of-unity Apostol-Bernoulli generating form, now obtained as one channel of an exact refinement bank. The Lerch representation uses the conventional analytic continuation described in~\cite{DLMF25}.

\FlowChapter{Bernoulli and Euler as low-pass and high-pass outputs}

Take $L=2$. The two children are
\[
2^{n-1}B_n\left(\frac a2\right)
\quad\text{and}\quad
2^{n-1}B_n\left(\frac{a+1}{2}\right).
\]

\begin{theorem}[Two-channel identity]
\[
\boxed{
W_{2,0}(n;a)=\frac{B_n(a)}{\sqrt2}
}
\]
and, for $n\ge1$,
\[
\boxed{
W_{2,1}(n;a)
=-\frac{n}{2\sqrt2}E_{n-1}(a).
}
\]
\end{theorem}

\begin{proof}
The first formula is Raabe multiplication. The second follows from
\[
E_{n-1}(a)
=\frac{2^n}{n}
\left[
B_n\left(\frac{a+1}{2}\right)
-B_n\left(\frac a2\right)
\right].
\]
\end{proof}

Thus Bernoulli and Euler polynomials are the DC and Nyquist channels of one exact Bernoulli refinement.

\FigureWide{11_bernoulli_euler_channels}{Two Bernoulli children transformed into Bernoulli and Euler channels.}{The transform changes coordinates without changing the information carried by the pair.}

\FlowChapter{Packet Wronskians and zero diagnostics}

Define
\[
\boxed{
\mathcal D_{j,\ell}(z;a)
=
\Wcal_j(z;a)\Wcal_\ell(z-1;a)
-
\Wcal_\ell(z;a)\Wcal_j(z-1;a).
}
\]

\begin{proposition}[Ratio derivative]
Where $\Wcal_\ell(z;a)\ne0$,
\[
\boxed{
\frac{\partial}{\partial a}
\left(
\frac{\Wcal_j(z;a)}{\Wcal_\ell(z;a)}
\right)
=
-\frac{z\mathcal D_{j,\ell}(z;a)}{\Wcal_\ell(z;a)^2}.
}
\]
\end{proposition}

\begin{proof}
Differentiate the ratio and use the order-lowering differential ladder for both channels.
\end{proof}

For the two-channel integer family, define
\[
\boxed{
\Wcal_n(a)
=
 nB_{n-1}(a)E_{n-1}(a)
-(n-1)B_n(a)E_{n-2}(a).
}
\]
This is four times the packet Wronskian of the Bernoulli and Euler channels at consecutive orders.

\begin{example}[Orders two through four]
\[
\Wcal_2(a)=\frac{3a^2-3a+1}{3},
\]
which has no real zeros. Also,
\[
\Wcal_3(a)=a^2(a-1)^2,
\]
while
\[
\Wcal_4(a)=\frac{a(a-1)}5
\left(5a^4-10a^3+5a-2\right).
\]
The fourth-order polynomial has nonreal roots.
\end{example}

\FigureWide{12_packet_wronskians}{Four low-order packet Wronskians.}{The zeros can be read directly from the ratio derivative and provide a concrete interlacing diagnostic.}
\FigureWide{13_wronskian_root_counts}{Real and nonreal roots for the first fifteen packet Wronskians.}{The finite computation rejects the naive conjecture that this entire Wronskian family is real-rooted.}

\begin{boundary}[What the negative control means]
The failure of universal real-rootedness does not remove the Wronskian route. It shows that a successful theorem must impose the right channel, order, parameter, or completed-function normalization. A blanket real-rootedness statement is false.
\end{boundary}

\FlowChapter{A worked cyclotomic norm over the third roots of unity}

Use the unnormalized $L=3$ detail
\[
D_1(n)=
\sum_{r=0}^{2}\omega_3^{-r}
3^{n-1}B_n\left(\frac r3\right).
\]
Write $D_1=A_n+B_n\omega_3$. Its field norm is
\[
N_n=A_n^2-A_nB_n+B_n^2.
\]

\begin{theorem}[Parity law for the third-root detail norm]
For every even $n\ge2$,
\[
D_1(n)=3^{n-1}\left[B_n-B_n\left(\frac13\right)\right]\in\Q
\]
and
\[
\boxed{N_n=\left\{3^{n-1}\left[B_n-B_n\left(\frac13\right)\right]\right\}^2.}
\]
For every odd $n\ge3$,
\[
D_1(n)=(\omega_3^2-\omega_3)3^{n-1}B_n\left(\frac13\right)
\]
and
\[
\boxed{N_n=3\left[3^{n-1}B_n\left(\frac13\right)\right]^2.}
\]
\end{theorem}

\begin{proof}
Bernoulli reflection gives $B_n(1-x)=(-1)^nB_n(x)$. For even $n$, the two child values at $1/3$ and $2/3$ agree, and $\omega_3+\omega_3^2=-1$ reduces the detail to the displayed rational number. For odd $n\ge3$, one has $B_n=0$ and $B_n(2/3)=-B_n(1/3)$. Since $(\omega_3^2-\omega_3)(\omega_3-\omega_3^2)=3$, the field norm is three times the indicated rational square.
\end{proof}

\begin{longtable}{rllll}
\caption{Exact first non-DC cyclotomic detail data for $L=3$ and $a=0$.}\label{tab:cyclotomic-norms}\\
\toprule
$n$ & $A_n$ & $B_n$ & $N_n=A_n^2-A_nB_n+B_n^2$ & factor pattern\\
\midrule
\endfirsthead
\toprule
$n$ & $A_n$ & $B_n$ & $N_n$ & factor pattern\\
\midrule
\endhead
1 & $-1/3$ & $1/3$ & $1/3$ & first-order exception\\
2 & $2/3$ & $0$ & $4/9$ & $9N_n$ is a square\\
3 & $-1/3$ & $-2/3$ & $1/3$ & $3N_n$ is a square\\
4 & $-4/3$ & $0$ & $16/9$ & $9N_n$ is a square\\
5 & $5/3$ & $10/3$ & $25/3$ & $3N_n$ is a square\\
6 & $26/3$ & $0$ & $676/9$ & $9N_n$ is a square\\
7 & $-49/3$ & $-98/3$ & $2401/3$ & $3N_n$ is a square\\
8 & $-328/3$ & $0$ & $107584/9$ & $9N_n$ is a square\\
9 & $809/3$ & $1618/3$ & $654481/3$ & $3N_n$ is a square\\
10 & $6710/3$ & $0$ & $45024100/9$ & $9N_n$ is a square\\
11 & $-20317/3$ & $-40634/3$ & $412780489/3$ & $3N_n$ is a square\\
12 & $-201772/3$ & $0$ & $40711939984/9$ & $9N_n$ is a square\\
\bottomrule
\end{longtable}


\FigureWide{14_cyclotomic_norms}{Growth of the exact $L=3$ cyclotomic detail norm.}{The parity-dependent square pattern follows from Bernoulli reflection symmetry.}

\InlinePart{M\"obius-Bernoulli transfer functions}

\FlowChapter{A generalized reciprocal L-function transfer family}

Let $\chi$ be a Dirichlet character. For parameters $a>1$, $\beta>0$, integer $k\ge0$, and real $t$, define
\[
\boxed{
\Ccal_{a,\beta,\chi}(k,t)
=
\sum_{j=0}^{k}
(-1)^j\binom{k}{j}
\frac1{L(a+\beta j+\ii t,\chi)}.
}
\]

\begin{theorem}[M\"obius form]
In the half-plane of absolute convergence,
\[
\boxed{
\Ccal_{a,\beta,\chi}(k,t)
=
\sum_{n=1}^{\infty}
\frac{\mu(n)\chi(n)}{n^{a+\ii t}}
\left(1-\frac1{n^\beta}\right)^k.
}
\]
\end{theorem}

\begin{proof}
Use
\[
\frac1{L(s,\chi)}
=
\sum_{n\ge1}\frac{\mu(n)\chi(n)}{n^s}
\]
and interchange the finite and absolutely convergent sums. The binomial sum becomes $(1-n^{-\beta})^k$.
\end{proof}

\begin{proposition}[Parameter-shift recursion]
\[
\boxed{
\Ccal_{a,\beta,\chi}(k+1,t)
=
\Ccal_{a,\beta,\chi}(k,t)
-
\Ccal_{a+\beta,\beta,\chi}(k,t).
}
\]
\end{proposition}

\begin{proof}
Apply Pascal's identity to the binomial coefficient and shift the index in the second sum.
\end{proof}

\FlowChapter{Exponential generating and Mellin transforms}

Define
\[
\boxed{
\Rcal_{a,\beta,\chi}(x,t)
=
x\sum_{n=1}^{\infty}
\frac{\mu(n)\chi(n)}{n^{a+\ii t}}
\e^{-x/n^\beta}.
}
\]

\begin{theorem}[Poissonization]
\[
\boxed{
\sum_{k=0}^{\infty}
\Ccal_{a,\beta,\chi}(k,t)
\frac{x^k}{k!}
=
\frac{\e^x}{x}\Rcal_{a,\beta,\chi}(x,t).
}
\]
\end{theorem}

\begin{proof}
For each integer $n$,
\[
\sum_{k\ge0}
\left(1-\frac1{n^\beta}\right)^k
\frac{x^k}{k!}
=
\exp\left[x\left(1-\frac1{n^\beta}\right)\right].
\]
Sum over $n$.
\end{proof}

\begin{theorem}[Mellin transform]
For absolute convergence, assume
\[
-1<\Re s<\frac{a-1}{\beta}-1.
\]
The resulting identity extends meromorphically wherever the reciprocal $L$-function side is defined.
\[
\boxed{
\int_0^\infty
\Rcal_{a,\beta,\chi}(x,t)x^{s-1}\dd x
=
\frac{\Gamma(s+1)}
{L(a+\ii t-\beta(s+1),\chi)}.
}
\]
\end{theorem}

\begin{proof}
Integrate term by term. The change of variable $x=n^\beta y$ gives
\[
\int_0^\infty x^s\e^{-x/n^\beta}\dd x
=n^{\beta(s+1)}\Gamma(s+1).
\]
The remaining Dirichlet series is the reciprocal $L$-function.
\end{proof}

\FlowChapter{The zeta specialization and the Bernoulli seed}

Take the principal character, $a=2$, $\beta=2$, and $t=0$. Then
\[
\boxed{
 c_k
=
\Ccal_{2,2}(k,0)
=
\sum_{j=0}^{k}
(-1)^j\binom{k}{j}
\frac1{\zeta(2j+2)}.
}
\]
The even zeta values satisfy
\[
\frac1{\zeta(2j+2)}
=
\frac{2(2j+2)!}
{(2\pi)^{2j+2}\abs{B_{2j+2}}}.
\]
If
\[
q=(2\pi)^{-2},
\qquad
A_j=\frac{2(2j+2)!}{\abs{B_{2j+2}}},
\]
then
\[
\boxed{
c_k=q(I-qE)^kA_0.
}
\]
The B\'aez-Duarte theorem~\cite{BaezDuarte2005,CisloWolf2008} states that RH is equivalent to
\[
\boxed{
 c_k=O_\varepsilon(k^{-3/4+\varepsilon})
\quad\text{for every }\varepsilon>0.
}
\]
Thus RH is a stability statement for repeated weighted finite differences of one normalized Bernoulli seed.

\FlowChapter{Scale localization and the Mertens wavelet}

Write
\[
f_k(x)=x^{-2}(1-x^{-2})^k.
\]
Then
\[
c_k=\sum_{n\ge1}\mu(n)f_k(n).
\]
The maximum of $f_k$ occurs at
\[
x=\sqrt{k+1}.
\]
More generally, the kernel
\[
x^{-a}(1-x^{-\beta})^k
\]
is centered at
\[
\boxed{
 x_*(k)=\left(1+\frac{\beta k}{a}\right)^{1/\beta}.
}
\]

\FigureWide{15_mobius_kernel_localization}{Rescaled principal transfer kernels.}{The curves collapse near one common profile after the coordinate is divided by $\sqrt{k}$.}

Let
\[
M(x)=\sum_{n\le x}\mu(n).
\]
Discrete summation by parts gives
\[
\boxed{
 c_k=
\sum_{n=1}^{\infty}
M(n)\left[f_k(n)-f_k(n+1)\right].
}
\]
The continuous wavelet is
\[
\boxed{
\psi_k(x)=-f_k'(x)
=2x^{-5}(1-x^{-2})^{k-1}(x^2-k-1).
}
\]
It has one sign change at $x=\sqrt{k+1}$ and zero integral.

\FigureWide{16_mertens_wavelet}{The Mertens wavelet at three scales.}{The sign change remains fixed in the normalized coordinate, making the scale selection explicit.}

\begin{theorem}[Mellin response of the wavelet]
For $k\ge1$ and $\Re s<3$,
\[
\boxed{
\int_1^\infty
\psi_k(x)x^{s-1}\dd x
=
\frac{s-1}{2}
\mathrm B\left(\frac{3-s}{2},k+1\right).
}
\]
\end{theorem}

\begin{proof}
First compute
\[
\int_1^\infty f_k(x)x^{s-1}\dd x
=\frac12\mathrm B\left(1-\frac s2,k+1\right)
\]
by the substitution $y=x^{-2}$. Integration by parts transfers the derivative from $f_k$ to $x^{s-1}$.
\end{proof}

Mellin inversion gives the exact N\"orlund-Rice form
\[
\boxed{
 c_k
=
\frac1{4\pi\ii}
\int_{\sigma-\ii\infty}^{\sigma+\ii\infty}
\frac{
\mathrm B(1-s/2,k+1)
}{\zeta(s)}\dd s,
\qquad 1<\sigma<2.
}
\]
A simple zero $\rho$ contributes asymptotically
\[
\boxed{
\frac{\Gamma(1-\rho/2)}{2\zeta'(\rho)}
 k^{\rho/2-1}.
}
\]
In logarithmic time $u=\log k$, an off-critical zero creates an exponentially growing normalized mode.

\FlowChapter{Diagonal packet energy}

Define
\[
E_{a,\beta,\chi}(k)
=
\sum_{n=1}^{\infty}
\frac{\mu(n)^2\abs{\chi(n)}^2}{n^{2a}}
\left(1-\frac1{n^\beta}\right)^{2k}.
\]
For the asymptotic theorem below, assume that $\chi$ is primitive with conductor $q$, and set
\[
\delta_\chi
=
\frac1{\zeta(2)}
\prod_{p\mid q}\frac{p}{p+1}.
\]

\begin{theorem}[Packet-energy asymptotic]
For $a>1/2$,
\[
\boxed{
E_{a,\beta,\chi}(k)
\sim
\frac{\delta_\chi}{\beta}
\Gamma\left(\frac{2a-1}{\beta}\right)
(2k)^{-(2a-1)/\beta}.
}
\]
\end{theorem}

\begin{proof}
The summatory density of squarefree integers coprime to $q$ is
\[
\sum_{n\le x}\mu(n)^2\abs{\chi(n)}^2
=\delta_\chi x+O_q(\sqrt x).
\]
Use partial summation. The kernel is concentrated at $n\asymp k^{1/\beta}$. On this scale,
\[
\left(1-n^{-\beta}\right)^{2k}
=\exp(-2kn^{-\beta})(1+o(1)).
\]
After the scaling $n=(2k)^{1/\beta}y$, the main term becomes
\[
\delta_\chi(2k)^{-(2a-1)/\beta}
\int_0^\infty y^{-2a}\e^{-y^{-\beta}}\dd y.
\]
The substitution $v=y^{-\beta}$ evaluates the integral as
\[
\frac1\beta\Gamma\left(\frac{2a-1}{\beta}\right).
\]
The square-root error in the summatory function is lower order after the same localization.
\end{proof}

For $a=\beta=2$,
\[
\boxed{
E_k
\sim
\frac{3}{2^{5/2}\pi^{3/2}}k^{-3/2}.
}
\]

\FigureWide{17_energy_asymptotic}{The scaled energy $k^{3/2}E_k$.}{The numerical curve approaches the exact constant predicted by the theorem.}

Define also the active squarefree mass
\[
S_k=
\sum_{n\ge1}
\frac{\mu(n)^2}{n^2}
\left(1-\frac1{n^2}\right)^k.
\]
The same argument gives
\[
S_k\sim\frac{3}{\pi^{3/2}}k^{-1/2}.
\]
Therefore
\[
\boxed{
\frac{S_k^2}{E_k}
\sim
\frac{3\cdot2^{5/2}}{\pi^{3/2}}\sqrt{k}.
}
\]
This is the effective number of active packets.

\FigureWide{18_mass_and_dimension}{Active squarefree mass and effective dimension.}{The number of packets contributing at one scale grows like a constant multiple of $\sqrt{k}$.}

\FlowChapter{M\"obius packet coherence and RH}

Define
\[
\boxed{
\Gamma_k=\frac{\abs{c_k}^2}{E_k}.
}
\]

\begin{theorem}[Coherence form of the B\'aez-Duarte criterion]
The Riemann hypothesis is equivalent to
\[
\boxed{
\Gamma_k=O_\varepsilon(k^\varepsilon)
\quad\text{for every }\varepsilon>0.
}
\]
\end{theorem}

\begin{proof}
The energy theorem gives $E_k\asymp k^{-3/2}$. If RH holds, the B\'aez-Duarte criterion gives $|c_k|\ll_\varepsilon k^{-3/4+\varepsilon}$. Hence $\Gamma_k\ll_\varepsilon k^{2\varepsilon}$, which is subpolynomial after renaming $2\varepsilon$. Conversely, subpolynomial coherence gives
\[
|c_k|^2\ll_\varepsilon k^{-3/2+\varepsilon},
\]
so $|c_k|\ll_\varepsilon k^{-3/4+\varepsilon}$ after another harmless renaming.
\end{proof}

The theorem separates the problem into two parts. The diagonal energy decay is already proved. The remaining RH content is the off-diagonal arithmetic coherence of the M\"obius signs.

\begin{table}[htbp]
\centering
\small
\caption{Selected numerical values of the principal Möbius-Bernoulli transfer.}
\label{tab:mobius-samples}
\begin{tabular}{rrrrr}
\toprule
$k$ & $c_k$ & $E_k$ & $\Gamma_k$ & $S_k^2/E_k$\\
\midrule
10 & -6.91391e-02 & 3.04674e-03 & 1.56896e+00 & 1.06483e+01\\
50 & -5.36535e-03 & 2.76532e-04 & 1.04100e-01 & 2.17944e+01\\
102 & -1.42289e-03 & 9.13011e-05 & 2.21752e-02 & 3.14049e+01\\
517 & -5.96831e-05 & 8.03172e-06 & 4.43501e-04 & 6.99561e+01\\
950 & -1.84349e-05 & 3.26717e-06 & 1.04018e-04 & 9.34893e+01\\
4800 & -8.04548e-07 & 2.86464e-07 & 2.25961e-06 & 2.10639e+02\\
9750 & -1.33337e-07 & 9.91180e-08 & 1.79370e-07 & 3.00488e+02\\
49238 & 1.69527e-08 & 8.73393e-09 & 3.29055e-08 & 6.73612e+02\\
100000 & 1.82845e-09 & 3.00705e-09 & 1.11179e-09 & 9.61804e+02\\
\bottomrule
\end{tabular}
\end{table}


\FigureWide{19_coherence_ratio}{The principal transfer amplitude, diagonal energy scale, and coherence ratio.}{The numerical values are finite diagnostics only. They do not establish the required asymptotic bound.}
\FigureWide{20_log_time_signal}{The normalized sequence $k^{3/4}c_k$ in logarithmic time.}{Zeros off the critical line would appear as expanding modes in this coordinate.}

\begin{boundary}[Why a generic frame estimate cannot prove RH]
If every active packet had the same sign, the coherent-to-diagonal ratio would grow like $\sqrt{k}$. The missing amplitude factor is $k^{1/4}$. Therefore Cauchy-Schwarz, Parseval, positivity, or a sign-free frame bound cannot prove RH. Any successful estimate must use arithmetic structure of $\mu(n)$.
\end{boundary}

\FlowChapter{Long-time mean power and a quantum correlation}

Define
\[
F_k(t)
=
\sum_{n=1}^{\infty}
\frac{\mu(n)}{n^{2+\ii t}}
\left(1-\frac1{n^2}\right)^k.
\]
Then
\[
F_k(0)=c_k.
\]

\begin{theorem}[Long-time diagonalization]
\[
\boxed{
\lim_{T\to\infty}
\frac1{2T}
\int_{-T}^{T}\abs{F_k(t)}^2\dd t
=E_k.
}
\]
\end{theorem}

\begin{proof}
The frequencies $\log n$ are distinct. For a finite truncation, every off-diagonal time average tends to zero. The coefficient sequence is square summable, so the truncation tail is controlled in mean square. Pass to the limit.
\end{proof}

Therefore
\[
\boxed{
\Gamma_k
=
\frac{\abs{F_k(0)}^2}
{\displaystyle\lim_{T\to\infty}\frac1{2T}\int_{-T}^{T}\abs{F_k(t)}^2\dd t}.
}
\]
RH says that the zero-time recurrence is only subpolynomially larger than the long-time mean power.

\FigureWide{21_long_time_power}{One finite M\"obius-Bernoulli correlation and its running mean power.}{The dashed line is the exact diagonal energy of the truncated packet family.}

For $a>1$ and $k\ge1$, work in $\ell^2(\{2,3,\ldots\})$. Define
\[
|\phi_a\rangle=\sum_{n\ge2}n^{-a/2}|n\rangle,
\]
\[
H|n\rangle=\log n|n\rangle,
\]
\[
K_\beta|n\rangle=-\log(1-n^{-\beta})|n\rangle,
\]
and
\[
\Mcal_\chi|n\rangle=\mu(n)\chi(n)|n\rangle.
\]
Then
\[
\boxed{
\Ccal_{a,\beta,\chi}(k,t)
=\langle\phi_a|
\Mcal_\chi
\e^{-kK_\beta}
\e^{-\ii tH}
|\phi_a\rangle.
}
\]
The logarithmic Hamiltonian $H$ is the same spectral object used in the recent zeta-governed quantum-dynamics construction of Wei and collaborators~\cite{WeiEtAl2026}. That work builds thermal weights and phases from the spectrum $\log n$. The new ingredients here are the M\"obius-character observable and the scale filter $\e^{-kK_\beta}$. A complete gate construction would require a reversible M\"obius or squarefree-phase oracle, which remains an engineering and complexity question.

\FigureWide{24_quantum_operator_pipeline}{Operator realization of the generalized transfer function.}{Scale evolution and real-time logarithmic evolution commute because both are diagonal in the integer basis.}

\FlowChapter{Arithmetic polyphase and character packets}

For a dyadic scale shell $I_J=[2^J,2^{J+1})$ and a residue class $r\pmod M$, define
\[
\boxed{
F_{k;J,M,r}(t)
=
\sum_{\substack{n\in I_J\\n\equiv r\pmod M}}
\frac{\mu(n)}{n^{2+\ii t}}
\left(1-\frac1{n^2}\right)^k.
}
\]
Then
\[
F_k(t)=\sum_J\sum_{r=0}^{M-1}F_{k;J,M,r}(t).
\]
Arithmetic refinement is exact:
\[
\boxed{
F_{k;J,M,r}
=
\sum_{\ell=0}^{L-1}
F_{k;J,LM,r+\ell M}.
}
\]
Scale refinement is exact because each shell splits into two half-shells. The diagonal energy also separates exactly over every complete packet family.

\FigureWide{22_polyphase_mobius_packets}{Residue packets of the principal M\"obius transfer at three scales.}{Positive and negative packets shift with $k$, while their exact sum remains $c_k$.}

For a Dirichlet character $\chi$, the character transfer
\[
\Ccal_{a,\beta,\chi}(k,t)
=
\sum_{n\ge1}
\frac{\mu(n)\chi(n)}{n^{a+\ii t}}
(1-n^{-\beta})^k
\]
provides a multiplicative-frequency coordinate. Known Riesz-type theorems give GRH criteria for continuous exponential kernels and related $L$-function classes~\cite{GargMaji2022,GuptaVatwani2022}. The corresponding discrete coherence criterion is a natural target, but no character-uniform theorem is claimed here.

\FigureWide{23_character_packets}{Principal and quadratic-character transfer amplitudes.}{The character coordinate can strongly alter cancellation even when the same scale filter is used.}

\InlinePart{Worked computational atlas and research program}

\FlowChapter{Worked example: complete three-channel exponential lift}

Start from
\[
e^z=\sum_{n\ge0}\frac{z^n}{n!}.
\]
Separate the exponents modulo three:
\[
\begin{aligned}
E_{3,0}(u)&=1+\frac{u}{3!}+\frac{u^2}{6!}+\cdots,\\
E_{3,1}(u)&=1+\frac{u}{4!}+\frac{u^2}{7!}+\cdots,\\
E_{3,2}(u)&=\frac1{2!}+\frac{u}{5!}+\frac{u^2}{8!}+\cdots.
\end{aligned}
\]
The packet system is
\[
\Theta_3
\begin{pmatrix}E_{3,0}\\E_{3,1}\\E_{3,2}\end{pmatrix}
=
\begin{pmatrix}
0&0&u\\
1&0&0\\
0&1&0
\end{pmatrix}
\begin{pmatrix}E_{3,0}\\E_{3,1}\\E_{3,2}\end{pmatrix}.
\]
Writing the rows explicitly gives
\[
\begin{aligned}
3uE_{3,0}'&=uE_{3,2},\\
(3u\partial_u+1)E_{3,1}&=E_{3,0},\\
(3u\partial_u+2)E_{3,2}&=E_{3,1}.
\end{aligned}
\]
Reconstruction yields
\[
 e^z=E_{3,0}(z^3)+zE_{3,1}(z^3)+z^2E_{3,2}(z^3).
\]
Every equation can be checked coefficient by coefficient. The matrix system is not a numerical approximation. It is the scalar equation in an exact residue coordinate system.

\FlowChapter{Worked example: Bernoulli packet refinement}

Take $a=1/3$, $M=2$, and refine by $L=2$. The parent packet is
\[
\Bfrak_{2,1}(z;1/3)
=-z2^{z-1}\zeta\left(1-z,\frac23\right).
\]
Its two children are
\[
\Bfrak_{4,1}(z;1/3)
=-z4^{z-1}\zeta\left(1-z,\frac13\right)
\]
and
\[
\Bfrak_{4,3}(z;1/3)
=-z4^{z-1}\zeta\left(1-z,\frac56\right).
\]
The multiplication formula at modulus two gives
\[
\zeta\left(1-z,\frac13\right)
+
\zeta\left(1-z,\frac56\right)
=2^{1-z}\zeta\left(1-z,\frac23\right).
\]
Multiplication by $-z4^{z-1}$ gives
\[
\boxed{
\Bfrak_{2,1}(z;1/3)
=
\Bfrak_{4,1}(z;1/3)
+
\Bfrak_{4,3}(z;1/3).
}
\]
At $z=6$, this becomes an exact identity among rational Bernoulli-polynomial values.

\FlowChapter{Worked example: third-root cyclotomic norm}

For $n=8$, the three integer child packets at $a=0$ are
\[
-\frac{729}{10},
\qquad
\frac{1093}{30},
\qquad
\frac{1093}{30}.
\]
Let $\omega^2+\omega+1=0$. The first non-DC detail is
\[
D_1
=-\frac{729}{10}
+\omega^2\frac{1093}{30}
+\omega\frac{1093}{30}
=-\frac{328}{3}.
\]
Its cyclotomic norm is
\[
\boxed{
N_8=D_1\overline{D_1}=\frac{107584}{9}=\left(\frac{328}{3}\right)^2.
}
\]
For odd $n\ge3$, the same calculation produces three times a rational square. The table in \cref{tab:cyclotomic-norms} records the first values.

\FlowChapter{Worked example: the principal transfer at order 1200}

For $k=1200$, define
\[
w_n=\frac{\mu(n)}{n^2}\left(1-\frac1{n^2}\right)^{1200}.
\]
The finite computation to $n=200000$ gives a zero-time amplitude of approximately
\[
\abs{F_{1200}(0)}=1.1153\times10^{-5}.
\]
The diagonal packet power is approximately
\[
E_{1200}=2.3091\times10^{-6}.
\]
Hence
\[
\Gamma_{1200}\approx5.3868\times10^{-5}.
\]
This is one finite numerical value. It is not an asymptotic theorem. Its purpose is to show how the exact definitions become a measurable signal-to-energy ratio.

\FlowChapter{Algorithmic workflow}

\begin{algorithm}[Generating a polyphase special-function system]
Given a scalar power series or differential equation:
\begin{enumerate}[label=\arabic*.]
\item Choose a modulus $M$ and separate Taylor exponents by residue.
\item Form $\mathbf F_M(u)$.
\item Replace multiplication by $z$ with $C_M(u)$.
\item Replace $z\frac{\dd}{\dd z}$ with $\Theta_M$.
\item Preserve operator order and derive the matrix system.
\item If a parent-child identity is available, refine packet addresses.
\item Apply a unitary transform to create low-pass and detail channels.
\item Derive inverse synthesis and energy identities.
\item Build Wronskians, recurrences, arithmetic traces, or coherence ratios suited to the question.
\item Test simple conjectures against exact and high-precision counterexamples before promoting them.
\end{enumerate}
\end{algorithm}

\FigureWide{25_function_family_map}{The generated function families and their main transfer routes.}{The same packet operations produce differential systems, refinable Bernoulli packets, cyclotomic details, Wronskians, and arithmetic coherence signals.}

\FlowChapter{Open problems generated by the calculus}

\section{Near-term problems}
\begin{enumerate}[label=\textbf{P\arabic*.}]
\item Determine sharp denominator and local valuation laws for cyclotomic Bernoulli details at rational $a$.
\item Classify the orders and parameter intervals for which selected packet Wronskians have one sign.
\item Determine when Bernoulli and Euler channel zeros interlace after completion or weighting.
\item Develop dominant-pole asymptotics for reciprocal and logarithmic derivatives of non-DC detail functions.
\item Find minimal matrix differential systems after symmetry reduction of the general polyphase lift.
\item Identify transforms for which packet children form a positive or totally positive kernel.
\end{enumerate}

\section{Major-conjecture interfaces}
\begin{enumerate}[label=\textbf{M\arabic*.}]
\item Prove subpolynomial M\"obius packet coherence in the principal zeta case.
\item Establish a character-uniform version compatible with GRH criteria.
\item Construct positive canonical systems from packet Gram matrices.
\item Connect local $p$-adic packet estimates to Archimedean coherence through a quantitative global height theorem.
\end{enumerate}

\begin{conjecture}[M\"obius packet coherence]
For every $\varepsilon>0$,
\[
\boxed{
\left|
\sum_{n=1}^{\infty}
\frac{\mu(n)}{n^2}
\left(1-\frac1{n^2}\right)^k
\right|^2
\le
C_\varepsilon k^\varepsilon
\sum_{n=1}^{\infty}
\frac{\mu(n)^2}{n^4}
\left(1-\frac1{n^2}\right)^{2k}.
}
\]
\end{conjecture}

By the proved energy asymptotic and the B\'aez-Duarte criterion, this conjecture is equivalent to RH. It is included as the sharp theorem target, not as a claimed result.

\FigureWide{26_status_map}{Status map for the paper.}{The map separates proved identities, classical inputs, finite computations, open special-function questions, and major conjecture interfaces.}

\FlowChapter{Verification}

Every figure and numerical table was generated by the accompanying Python source. The symbolic suite verifies coefficient reconstruction, the Weyl commutator, hypergeometric exponential packets, Bernoulli and Euler channel identities, complex-order Bernoulli reconstruction and refinement, Lerch closed forms, and packet Wronskian formulas.

\begin{table}[htbp]
\centering
\small
\caption{Exact and high-precision verification summary.}
\label{tab:checks}
\begin{tabular}{lrr}
\toprule
Check family & checks & passed\\
\midrule
Bernoulli & 15 & 15\\
Euler & 9 & 9\\
Lerch & 3 & 3\\
Weyl & 4 & 4\\
exponential & 9 & 9\\
packet & 7 & 7\\
polynomial & 4 & 4\\
\midrule
Total & 51 & 51\\
\bottomrule
\end{tabular}
\end{table}


The asymptotic plots use a M\"obius sieve through $200000$. The largest plotted transfer order is $100000$. These computations illustrate the proved formulas and expose possible failure modes. They do not replace the asymptotic theorems or the open coherence estimate.

\FlowChapter{Conclusion}

The calculus generates a coherent special-function theory rather than a list of renamed scalar functions. The defining structure is:
\[
\boxed{
\text{analysis}
\longrightarrow
\text{matrix transfer}
\longrightarrow
\text{refinement}
\longrightarrow
\text{unitary detail transform}
\longrightarrow
\text{synthesis}
\longrightarrow
\text{certificate}.
}
\]

The general polyphase lift transports scalar differential equations into exact matrix systems. Bernoulli-Hurwitz packets extend Raabe multiplication to complex order. Their cyclotomic transform produces Lerch details with perfect reconstruction. The two-channel case unifies Bernoulli and Euler polynomials. Packet Wronskians create direct zero diagnostics and also reveal false broad conjectures. The M\"obius-Bernoulli transfer family turns reciprocal $L$-values into scale-localized signals whose diagonal energy is explicit. In the zeta specialization, RH is precisely the remaining coherence estimate.

The practical next step is not to introduce more notation. It is to prove one nontrivial property of a generated family that is not automatic from the scalar parent. The most accessible targets are denominator laws, Wronskian signs, symmetry reductions, and reciprocal-detail recurrences. The most difficult target is the M\"obius coherence inequality.

\appendix

\FlowChapter{Formula registry}

\begin{longtable}{L{0.25\textwidth}L{0.65\textwidth}}
\toprule
Object & Defining formula\\
\midrule
\endfirsthead
\toprule
Object & Defining formula\\
\midrule
\endhead
Polyphase section & $F_{M,r}(u)=\sum_{m\ge0}a_{Mm+r}u^m$\\
Carry shift & $C_M(u)_{0,M-1}=u$, $C_M(u)_{r,r-1}=1$\\
Packet Euler operator & $\Theta_M=Mu\partial_uI+\diag(0,\ldots,M-1)$\\
Bernoulli interpolation & $\Bfrak(z;a)=-z\zeta(1-z,a)$\\
Bernoulli packet & $\Bfrak_{M,r}(z;a)=-zM^{z-1}\zeta(1-z,(a+r)/M)$\\
Cyclotomic detail & $\Wcal_{M,r}^{(L,j)}=L^{-1/2}\sum_\ell\omega_L^{-j\ell}\Bfrak_{LM,r+\ell M}$\\
Cyclotomic EGF & $\sum_nW_{L,j}(n;a)t^n/n!=t\e^{at}/[\sqrt L(\omega_L^{-j}\e^t-1)]$\\
Packet Wronskian & $\mathcal D_{j,\ell}=\Wcal_j(z)\Wcal_\ell(z-1)-\Wcal_\ell(z)\Wcal_j(z-1)$\\
General transfer & $\Ccal_{a,\beta,\chi}(k,t)=\sum_{j=0}^k(-1)^j\binom{k}{j}/L(a+\beta j+\ii t,\chi)$\\
M\"obius form & $\Ccal=\sum_n\mu(n)\chi(n)n^{-a-\ii t}(1-n^{-\beta})^k$\\
Continuous companion & $\Rcal=x\sum_n\mu(n)\chi(n)n^{-a-\ii t}\e^{-x/n^\beta}$\\
Packet energy & $E_{a,\beta,\chi}=\sum_n\mu(n)^2|\chi(n)|^2n^{-2a}(1-n^{-\beta})^{2k}$\\
Coherence ratio & $\Gamma=|\Ccal|^2/E$\\
\bottomrule
\end{longtable}

\FlowChapter{Complete Python source}
The following program generates every figure, table, and numerical record used in this paper.

\lstinputlisting{projector_generated_special_functions_listing.py}

\FlowChapter{Bibliographic note}

\begin{thebibliography}{99}

\bibitem{DLMF24}
NIST Digital Library of Mathematical Functions,
Chapter 24, ``Bernoulli and Euler Polynomials,''
particularly Section 24.4 on difference, derivative, and multiplication formulas.

\bibitem{DLMF25}
NIST Digital Library of Mathematical Functions,
Sections 25.11 and 25.14 on the Hurwitz zeta function and Lerch transcendent.

\bibitem{BaezDuarte2005}
L. B\'aez-Duarte,
``A Sequential Riesz-like Criterion for the Riemann Hypothesis,''
\emph{International Journal of Mathematics and Mathematical Sciences},
2005, Article 141565, 3527--3537.

\bibitem{CisloWolf2008}
J. Cis\l o and M. Wolf,
``On the Riesz and B\'aez-Duarte Criteria for the Riemann Hypothesis,''
2008, arXiv:0807.2971.

\bibitem{GargMaji2022}
M. Garg and B. Maji,
``Hardy-Littlewood-Riesz Type Equivalent Criteria for the Generalized Riemann Hypothesis,''
2022, arXiv:2208.07596.

\bibitem{GuptaVatwani2022}
S. Gupta and A. Vatwani,
``Riesz Type Criteria for $L$-Functions in the Selberg Class,''
2022, arXiv:2211.02954.

\bibitem{WeiEtAl2026}
S. Wei, Y. Zhai, Q. Lu, W. Yang, P. Gao, C. Wei, J. Song, F. Nori, T. Xin, and G. Long,
``The Riemann Hypothesis Manifested in Dynamical Quantum Phase Transitions,''
\emph{Nature Communications}, accepted 18 June 2026,
DOI: 10.1038/s41467-026-74935-8.

\bibitem{KomorebiBernoulli}
Komorebi P.,
``Bernoulli Transfer-Block Reformulation of Zeta,''
unpublished proof-development note.

\end{thebibliography}

\end{document}
